\section{本讲小结}

干涉强度来自复振幅相加后的模平方。相位差决定明暗位置，振幅匹配与相干性决定对比度；不同频率在通常观测时间内各自形成图样后叠加强度。双缝计算应先定义有符号光程差，加入薄片引起的额外光程和反射相位，再使用明暗条件。薄膜的传播光程差为 $2n_ft\cos r$，相位相消还须与振幅匹配区别。

单缝衍射把缝内连续微元相加，多缝干涉再把不同缝的复振幅相加。缝宽决定包络，周期决定主极大候选位置，缝数决定细峰宽度；缺级由单缝零点与缝间相长条件共同确定。有限包络下，总强度峰位还可能偏离干涉因子的极大位置，计数与求实际峰位应分别处理。

偏振片保留电场在透振轴上的分量。自然光经第一片的强度减半，后续各片则按当前偏振方向逐次使用马吕斯定律。旋光改变偏振方向，吸收改变强度；反求浓度须固定波长、温度与溶液厚度，并明确所采用的检偏方法。

光电效应中，单个光子的能量由频率决定，光子到达率还取决于功率。逸出功规定阈值，截止电压测量最大初动能，饱和电流反映被收集的电子数。对同一材料作差可消去逸出功；分析电流方向时先沿装置追踪电子，再反向标出常规电流。结果应同时符合能量守恒、量纲以及图中的电场和电子受力方向。
